\section{Colaboración de ingeniería y flujo de cambios}
\subsection{Perspectiva de gestión de proyectos}
Los equipos que trabajan junto con un Agent necesitan redefinir el «estado de la tarea»: los patches que genera el Agent deben incorporarse a una lista de revisión pendiente y mantenerse alineados con el backlog, los issues y la descomposición de tareas. Graham recomienda usar el enfoque de «los humanos se encargan de los requisitos + el Agent se encarga de la ejecución» para dividir los requisitos en varias tareas pequeñas, y luego registrar la actividad del Agent en el sprint board.

\begin{itemize}
    \item Tratar la fase de planeación del Agent como una «tarea de desarrollo» revisada por el PM o el tech lead.
    \item En la fase Act del Agent, conservar un mecanismo de confirmación humana que garantice que solo se ejecuten los comandos aprobados con los permisos correspondientes.
    \item En la fase Verify, publicar el reporte de pruebas en el issue para que QA lo revise en la regresión.
\end{itemize}

\subsection{Entrega y comunicación de riesgos}
Cuando el Agent ayuda a implementar un patch, el estado de la entrega debe identificarse claramente como «borrador del Agent + confirmación humana». Si el Agent falla, hay que notificar automáticamente al owner correspondiente, con el registro y el diff adjuntos, para que pueda hacerse cargo de inmediato.

 \begin{knowledgebox}{Matriz de comunicación}
Registro rápido: el Agent falla → se notifica al owner; el Agent tiene éxito pero requiere revisión → se da seguimiento al PR; el Agent ejecuta un comando peligroso → se solicita aprobación.
\end{knowledgebox}

\begin{warningbox}{No tratar al Agent como una caja negra}
En la comunicación hay que explicar qué está haciendo el Agent y por qué lo hace, para reducir el temor del equipo a confiar en él y garantizar que, en caso de falla, alguien pueda tomar el relevo con rapidez.
\end{warningbox}

\subsection{Resumen de la sección}
\begin{itemize}
    \item La gestión de proyectos necesita tratar la salida del Agent como un elemento de trabajo de primera clase e incorporarlo al tablero.
    \item La comunicación de riesgos y el mecanismo de aprobación son requisitos previos para el despliegue del Agent a gran escala.
\end{itemize}

